\documentclass[11pt]{article}
\usepackage[a4paper,margin=25mm]{geometry}
\usepackage{amsmath,amssymb,amsthm,booktabs,array}
\usepackage[T1]{fontenc}
\usepackage[colorlinks=true,linkcolor=blue,citecolor=blue,urlcolor=blue]{hyperref}
\newtheorem{theorem}{Theorem}
\newtheorem{lemma}{Lemma}
\newtheorem{corollary}{Corollary}
\newcommand{\tr}{\operatorname{Tr}}
\newcommand{\ket}[1]{\lvert #1\rangle}
\newcommand{\bra}[1]{\langle #1\rvert}
\newcommand{\EG}{E_{\mathrm G}}
\newcommand{\eps}{\epsilon_{\mathrm o}}
\title{A strict gap in the noisy-GHZ equality question of Xu et al.,\protect\\
Section 4.2, page 12 of arXiv:2211.05563v1}
\author{Haobo Ma\thanks{ChronoAI Pte Ltd, Singapore.}\and
Wenlin Zhang\thanks{National University of Singapore; ChronoAI Pte Ltd, Singapore.}}
\date{November 2, 2026\quad\textnormal{CC BY 4.0}}
\begin{document}
\maketitle
\begin{abstract}
The question in Xu et al., Section 4.2, page 12 of arXiv:2211.05563v1,
asks whether the optimized witness bound equals the convex-roof geometric
measure of genuine multipartite entanglement for noisy GHZ states.
We give a strict separation for every party count $n\geq3$ and local dimension
$d\geq3$, at an explicitly specified noise probability.
For three qutrits at noise $p=3/8$, the optimized witness bound is exactly
$1/9$, whereas the convex roof is at least $13/108$.
The proof combines a biseparable overlap inequality with a positive-definite
matrix and identifies the extra information supplied by the population
outside the witness-support subspace.
\end{abstract}

\section{The question and the witness}
Xu et al.\ \cite[Section 4.2, Eqs.~(21)--(26), pp.~11--12]{xu}
ask whether $\eps^{n,d}(p)=\EG(\rho_{n,d}(p))$ for their optimized GHZ witness.
Here the geometric measure is relative to \emph{biseparable} pure states,
not just fully product states. The convex-roof convention and the
single-observable optimization are those in that source.
The optimization of entanglement bounds from witness data is developed
by G\"uhne, Reimpell and Werner\ \cite{guhne}; their subsequent work treats
incomplete information and multiple observables\ \cite[Section II]{incomplete}.
The fidelity comparison in
Xu et al., Eq.~(22), is based on Dai et al.\ \cite{dai}.
Our separating inequality uses the same witness as Xu et al., without
changing its normalization or its set of biseparable states.

Let $\mathcal H=(\mathbb C^d)^{\otimes n}$, $M=d^n$, and put
\[
 \ket g=\frac1{\sqrt d}\sum_{j=0}^{d-1}\ket j^{\otimes n},\qquad
 G=\ket g\bra g,\qquad \rho(p)=\frac pM I+(1-p)G.
\]
Let $P$ be the diagonal projector onto computational words containing at
most two distinct symbols, and let $Q=I-P$. Distinct arrangements in
Xu et al., Eq.~(23), are counted once; their witness is precisely
\begin{equation}\label{eq:W}
 W=P/d-G.
\end{equation}
There are $d$ constant words and $\binom d2(2^n-2)$ words containing exactly
two symbols. Thus
\begin{equation}\label{eq:rank}
 R=\tr P=d+\binom d2(2^n-2),\qquad
 w(p)=\tr(\rho(p)W)=-\frac{d-1}{d}+k p,
 \quad k=\frac{d-1}{d}+\frac{R/d-1}{M}.
\end{equation}
The complement has positive rank when $n,d\geq3$, since a word using three
different symbols then exists.
For $d=2$, every computational word lies in $P$, so $W=I/2-G$ is the
fidelity witness. The optimized witness and fidelity bounds therefore
coincide, as stated in the caption of Xu et al., Fig.~3.

For a unit vector $\psi$ and a density operator $\rho$, define
\[
 \EG(\psi)=1-\max_{\phi\text{ biseparable},\ \|\phi\|=1}
                   |\langle\phi,\psi\rangle|^2,\qquad
 \EG(\rho)=\inf_{\rho=\sum_i t_i\ket{\psi_i}\bra{\psi_i}}
                    \sum_i t_i\EG(\psi_i),
\]
where the ensembles are finite, the weights are nonnegative and sum to one,
and their vectors have unit norm. For real $r$, set
\[
 \widehat E(rW)=\sup_{\|\psi\|=1}\{r\langle\psi,W\psi\rangle-\EG(\psi)\},
 \qquad \epsilon(w)=\sup_{r\in\mathbb R}\{rw-\widehat E(rW)\}.
\]
The expression for $\widehat E$ agrees with the two-vector optimization
in Xu et al., Eq.~(25). For the noisy state, $\eps^{n,d}(p)=\epsilon(w(p))$.
Every affine lower bound on pure-state entanglement extends to the
convex roof by averaging, so $\epsilon(w(p))\leq\EG(\rho(p))$.

\begin{theorem}\label{thm:all}
For every $n,d\geq3$, define
\[
 T=d^2-3d+3,\quad
 w_*=-\frac{2d-3}{dT},\quad
 p_* =\frac{(d-2)^2}{Tk},\quad
 e_* =\frac1{(d-1)^3+1},\quad
 h=\frac1{d^2(d-2)},
\]
with $k$ as in \eqref{eq:rank}. Then $0<p_*<1$, $w(p_*)=w_*$, and
\begin{equation}\label{eq:gap}
 \eps^{n,d}(p_*)=e_* < e_*+h p_*\left(1-\frac RM\right)
                  \leq \EG(\rho(p_*)).
\end{equation}
In particular, for $n=d=3$ one has $p_*=3/8$ and
\begin{equation}\label{eq:qutrit}
 \eps^{3,3}(3/8)=\frac19=\frac{12}{108}
       <\frac{13}{108}\leq\EG(\rho_{3,3}(3/8)).
\end{equation}
\end{theorem}

\section{A two-observable lower bound}
The following overlap inequality is the reason that the same certificate
works in all dimensions and party counts.
\begin{lemma}\label{lem:cut}
For every biseparable unit vector $\phi$,
$d\langle\phi,G\phi\rangle\leq\langle\phi,P\phi\rangle$.
\end{lemma}
\begin{proof}
Choose a nonempty proper cut $A\mid A^c$ across which
$\phi=u\otimes v$. For each symbol $j$, let $u_j$ and $v_j$ be the
amplitudes of $u$ and $v$ on the constant words with symbol $j$.
By Cauchy--Schwarz,
\[
 d\langle\phi,G\phi\rangle
 =\left|\sum_j u_jv_j\right|^2
 \leq \left(\sum_i|u_i|^2\right)\left(\sum_j|v_j|^2\right)
 =\sum_{i,j}|u_iv_j|^2
 \leq\langle\phi,P\phi\rangle.
\]
The last sum is the mass on the $d^2$ distinct words constant on each
side of this cut. They all use at most two symbols and hence lie in $P$.
This argument applies to every possible cut.
\end{proof}

Put
\begin{equation}\label{eq:parameters}
 a=\frac{d-1}{d},\quad b=\frac{(d-1)^2}{d(d-2)},\quad
 s=\frac{d-1}{d(d-2)}=b-a,\quad
 D=aG+b(P-G)+Q.
\end{equation}
The three summands are supported on mutually orthogonal subspaces.
Thus $D$ is positive definite and
\[
 D^{-1}=\frac d{d-1}G+\frac{d(d-2)}{(d-1)^2}(P-G)+Q.
\]
Writing $x=\langle\phi,G\phi\rangle$ and
$y=\langle\phi,P\phi\rangle$, Lemma~\ref{lem:cut} gives
\begin{align*}
 \langle\phi,D^{-1}\phi\rangle
 &=\frac{d(d-2)}{(d-1)^2}y+\frac d{(d-1)^2}x+(1-y)\\
 &\leq\frac{d(d-2)+1}{(d-1)^2}y+(1-y)=1.
\end{align*}
For any unit $\psi$, weighted Cauchy--Schwarz therefore gives
\[
 |\langle\phi,\psi\rangle|^2
 \leq\langle\phi,D^{-1}\phi\rangle\langle\psi,D\psi\rangle
 \leq\langle\psi,D\psi\rangle.
\]
Maximizing over biseparable $\phi$ yields
$\EG(\psi)\geq1-\langle\psi,D\psi\rangle$. Direct substitution in
\eqref{eq:parameters} gives $D=(1+h)I+sW-hQ$. Averaging over any
pure-state ensemble proves the useful bound
\begin{equation}\label{eq:certificate}
 \boxed{\ \EG(\rho)\geq-s\tr(\rho W)+h\tr(\rho Q)-h.\ }
\end{equation}
This is an analytic inequality for every density operator on $\mathcal H$.
It displays information that is absent when only $\tr(\rho W)$ is specified.

\section{An exact witness value and strict separation}
Let $m=d-1$, $N=m^4+m$, and consider
\[
 \ket\chi=\frac1{\sqrt N}\left(m^2\ket0^{\otimes n}
                          +\sum_{j=1}^{d-1}\ket j^{\otimes n}\right).
\]
Across every nontrivial cut, this is a Schmidt decomposition.
Its squared coefficients are $m^4/N$ and $1/N$, with the latter repeated
$d-1$ times. The maximal squared overlap
with a product vector across that cut is its largest squared Schmidt
coefficient: this follows by applying Cauchy--Schwarz to the diagonal
Schmidt matrix, whose operator norm is $m^2/\sqrt N$.
Maximizing over all cuts consequently gives
\begin{align*}
 \EG(\chi)&=\frac mN=e_*, & \langle\chi,P\chi\rangle&=1,\\
 \langle\chi,G\chi\rangle&=\frac{(m^2+m)^2}{dN}=\frac{d-1}{T}, &
 \langle\chi,W\chi\rangle&=w_*.
\end{align*}
The scalar identity $-s w_*-h=e_*$ now determines the optimized witness
bound exactly. Indeed, \eqref{eq:certificate} and $Q\succeq0$ imply
$\EG(\psi)\geq-s\langle\psi,W\psi\rangle-h$ for every pure unit vector.
Hence $\widehat E(-sW)\leq h$, and $\chi$ attains equality. Evaluating
the witness optimization at $r=-s$ gives $\epsilon(w_*)\geq e_*$.
Conversely, using $\chi$ in the defining supremum for $\widehat E(rW)$
gives $rw_*-\widehat E(rW)\leq\EG(\chi)=e_*$ for every real $r$.
Thus $\epsilon(w_*)=e_*$.

For completeness, $w_*$ lies strictly between $-(d-1)/d$ and $0$:
the positive difference from the left endpoint is $(d-2)^2/T$.
Also $k>0$ and $w(1)=(R/d-1)/M\geq0$.
Since $w(p)$ is affine and strictly increasing, its solution $w(p_*)=w_*$
lies in $(0,1)$ and has the value in Theorem~\ref{thm:all}.
Finally, $\tr(\rho(p_*)Q)=p_*(1-R/M)>0$. Substitution in
\eqref{eq:certificate} proves \eqref{eq:gap}.
For $n=d=3$, $R=21$, $M=27$, $k=8/9$, $s=2/3$, and $h=1/9$.
The additional term is $(1/9)(3/8)(6/27)=1/108$, proving
\eqref{eq:qutrit}. This completes the proof of Theorem~\ref{thm:all}.

\begin{corollary}\label{cor:interval}
For three qutrits, the inequality
$\eps^{3,3}(p)<\EG(\rho_{3,3}(p))$ holds throughout
$27/74<p<9/22$.
\end{corollary}
\begin{proof}
Equation~\eqref{eq:certificate} gives
$\EG(\rho(p))\geq L(p)=1/3-46p/81$.
A finite ensemble with witness mean $w$ and entanglement average $u$
always gives $\epsilon(w)\leq u$: each affine expression defining
$\epsilon$ is bounded by that average.
For $0\leq p\leq3/8$, mix $G$ with $\ket\chi\bra\chi$ using weight
$8p/3$ on the latter. Its witness mean is $w(p)$ and its entanglement
average is $U_-(p)=2/3-40p/27$, since $\EG(g)=2/3$ and
$\EG(\chi)=1/9$.
For $3/8\leq p\leq3/4$, mix $\ket\chi\bra\chi$ with
$\ket{000}\bra{000}$ using weight $2-8p/3$ on the former.
Its witness mean is again $w(p)$, and its average is
$U_+(p)=2/9-8p/27$.
Now $L-U_-=-1/3+74p/81>0$ for $p>27/74$, and
$L-U_+=1/9-22p/81>0$ for $p<9/22$.
The two ranges meet at $p=3/8$, proving the claim.
\end{proof}

\section{Source coverage and verification scope}
\begin{center}
\small
\begin{tabular}{@{}p{0.37\linewidth}p{0.57\linewidth}@{}}
\toprule
Source item in Xu et al.\ \cite{xu} & Coverage here\\
\midrule
Section 4.2, p.~12: noisy-GHZ equality & Refuted and replaced by
Theorem~\ref{thm:all}; Corollary~\ref{cor:interval} gives a noise interval.\\
Fig.~3: $d=3$, $n=3,5,7,9$ & Theorem~\ref{thm:all} gives a strict gap
at $p_*$ for each of these cases, and for every $n,d\geq3$.\\
Fig.~3 caption: qubit witness and fidelity bounds & Proved equal by
$P=I$ for $d=2$; this identifies the two bounds, not the convex roof.\\
Eqs.~(24)--(26): witness optimization & At the stated $w_*$, the value
is exactly $e_*$; the proof uses Eq.~(25)'s definition directly.\\
Section 3.3, Appendix B.3: witness form & Eq.~\eqref{eq:W} is the
equal-Schmidt-coefficient specialization of their witness; no optimality
claim about that witness is contradicted.\\
Section 3.2 and Appendix B.2: systematic graph-witness improvement & Not covered.\\
Section 3.4: broad noise-robustness conjecture; Section 4.1: general
unfaithfulness; Appendix B.4, Eq.~(76): singlet-witness conjecture & Not covered.\\
\bottomrule
\end{tabular}
\end{center}

The pinned Lean source
\href{https://github.com/the-omega-institute/trureturing/blob/cc684e6326490f9173bfefcadb0294ee6ccbb0b5/D5/S3/Quantum/Entanglement/Witness/NoisyGhzWitnessBoundGap.lean}{\texttt{NoisyGhzWitnessBoundGap.lean}}
contains the declaration
\begin{quote}\small\ttfamily
D5.S3.Quantum.Entanglement.Witness.\\
NoisyGhzWitnessBoundGap.result
\end{quote}
The source revision is
\texttt{cc684e6326490f9173bfefcadb0294ee6ccbb0b5}; the module path is
\nolinkurl{D5/S3/Quantum/Entanglement/Witness/NoisyGhzWitnessBoundGap.lean}.
The declaration's statement is
$\neg\,\texttt{claim}$. Here \texttt{claim} is equality for all
$n,d\geq3$ and all $0\leq p\leq1$.
Its proof combines \texttt{epsilon\_upper}, the bound
$\eps^{3,3}(3/8)\leq1/9$, with \texttt{roof\_lower}, the bound
$13/108\leq\EG(\rho_{3,3}(3/8))$.
The exact equality $\eps^{3,3}(3/8)=1/9$, the extension to every
$n,d\geq3$, and Corollary~\ref{cor:interval} are proved in writing above.
The accompanying \texttt{verify.py} checks the three-qutrit matrices
with rational arithmetic, enumerates selected finite support ranks, and
checks the stated scalar identities on explicitly reported finite ranges.
These finite calculations complement the written all-parameter proof.
The value of the convex roof itself is bounded, not determined exactly.

\begin{thebibliography}{9}
\bibitem{xu} X.-Y. Xu, Q. Zhou, S. Zhao, S.-M. Hu, L. Li, N.-L. Liu,
and K. Chen, \emph{A generic framework for genuine multipartite entanglement
detection}, \href{https://arxiv.org/abs/2211.05563v1}{arXiv:2211.05563v1}
(2022), Section 4.2 and Appendix B.3.
Journal version: \emph{General framework for genuine multipartite entanglement
detection}, Physical Review A \textbf{107}, 052405 (2023),
\href{https://doi.org/10.1103/PhysRevA.107.052405}{doi:10.1103/PhysRevA.107.052405}.
\bibitem{guhne} O. G\"uhne, M. Reimpell, and R. F. Werner,
\emph{Estimating entanglement measures in experiments},
Physical Review Letters \textbf{98}, 110502 (2007),
\href{https://doi.org/10.1103/PhysRevLett.98.110502}{doi:10.1103/PhysRevLett.98.110502}.
\bibitem{incomplete} O. G\"uhne, M. Reimpell, and R. F. Werner,
\emph{Lower bounds on entanglement measures from incomplete information},
\href{https://arxiv.org/abs/0802.1734}{arXiv:0802.1734} (2008), Section II.
\bibitem{dai} Y. Dai, Y. Dong, Z. Xu, W. You, C. Zhang, and O. G\"uhne,
\emph{Experimentally accessible lower bounds for genuine multipartite
entanglement and coherence measures}, Physical Review Applied \textbf{13},
054022 (2020),
\href{https://doi.org/10.1103/PhysRevApplied.13.054022}{doi:10.1103/PhysRevApplied.13.054022};
\href{https://arxiv.org/abs/1810.07492}{arXiv:1810.07492}.
\end{thebibliography}
\end{document}
